\subsection*{Fixed-window spectral ordering and ordinary ground overlap}

The fixed-window computer-assisted positivity certificate can be shifted to count its low
eigenvalues.  The thresholds in the two reflection sectors need not
coincide.  In particular, an even trial vector only sees the even spectral
gap when its overlap with the ground state is estimated.

\begin{proposition}[A fixed-window simple even ground state]
\label{prop:v117-ground-order}
Let $\mu_0\le\mu_1\le\cdots$ be the eigenvalues, counted with
multiplicity, of the complete canonical operator $\mathsf W_3$.
Then its ground state is simple and even, and
\begin{equation}\label{eq:v117-ground-order}
 0<\mu_0<3.644\cdot10^{-38},\qquad \mu_1>10^{-36}.
\end{equation}
Within the even sector alone, every eigenvalue other than $\mu_0$
exceeds $10^{-34}$.  These statements concern the complete operator,
including its infinite Fourier complement, at the single window
$\lambda=3$.
\end{proposition}
\begin{proof}[Shifted inertia and a matching trial direction]
For a threshold $a>0$ replace every diagonal $d_n$ by $d_n-a$ and
every tail lower weight $g_n$ by $g_n-a$.  Keep the same exact dyadic
witness $Z$ and the same head, intermediate and remote cutoffs as in
Proposition~\ref{prop:v116-window-positive}.  If $T_a=T-aI$, then
\[
 K_{Z,a}=K_{Z,0}-a(I+Z^*Z),\qquad
 R_a=B-T_aZ=R_0+aZ.
\]
Thus all intermediate residual rows change and are recomputed.  The
remote moment bound is unchanged, because those rows have no diagonal
entry and $G=\binom{I_{\rm head}}{-Z}$ is unchanged.  The diagonal inverse comparison gives
\[
 S_a:=F-aI-B^*T_a^{-1}B\succeq
 L_a:=K_{Z,a}-\sum_{N<n\le J}\frac{R_{a,n}^*R_{a,n}}{g_n-a}
                      -\frac{U_J}{g_{J+1}-a}.
\]
The sum uses normalized parity rows.  All denominators are strictly
positive by the earlier entire-tail bounds.

The signed interval LDL certificate for the even lower matrix at
$a_{\rm e}=10^{-34}$ has $257$ nonzero pivots, exactly one of which
is negative (zero-based index $24$).  For the odd lower matrix at
$a_{\rm o}=10^{-36}$ all $512$ pivots are strictly positive.
The complete signed elimination retains negative pivots with their
signs; it does not stop at the first negative pivot.  The exact dyadic
witnesses and every pivot enclosure are retained in the bundle;
\texttt{certify\_shifted\_inertia.py} recomputes all shifted residual rows.
The exact parameters are those of Proposition~\ref{prop:v116-window-positive}:
$(N,M,J,r)=(256,512,4096,80)$ for the even sector and
$(512,1024,4096,100)$ for the odd sector. Both certificates pass
at 768 bits and again at 896 bits using identical dyadic witnesses.
By finite-dimensional min--max (see also \cite[Section~4.3]{Teschl2009}), $S_{a_{\rm e}}$ has at most one
nonpositive eigenvalue, because the second eigenvalue of its lower
matrix is strictly positive.  The odd Schur matrix is positive definite.

A matching negative even direction is supplied by the existing
finite-supported candidate of
Proposition~\ref{prop:v112-finite-certificate}. More precisely, in that
$N=64$ certificate let $u$ be its normalized even vector, let $Q$ project
onto its even orthogonal complement. Set
$W_{64}=P_{64}\mathsf W_3P_{64}$, $\alpha=\ip{W_{64}u}{u}$,
$C=QW_{64}Q$ on that complement, and $r=QW_{64}u$. Write $A=C-\alpha I$, $z=A^{-1}r$, $q=\norm z$ and
$E=\ip zr$.  The vector
\[
 v=\frac{u-z}{\sqrt{1+q^2}}
\]
has norm one and belongs to the finite Fourier core.  Direct expansion
gives its exact complete-form Rayleigh value
\[
 \beta=QW_3(v,v)=\alpha-\frac{E}{1+q^2}.
\]
The previously certified enclosures imply
\[
 \alpha<5.312382\cdot10^{-38},\quad
 E>1.668982\cdot10^{-38},\quad q<0.000216,
 \qquad \beta<3.644\cdot10^{-38}.
\]
The last strict scalar comparison is separately checked by
\texttt{certify\_ground\_overlap\_constants.py}.  Because $v$ has
finite support, its finite-compression Rayleigh value equals its
complete-form value exactly.  No approximation of the omitted output
coefficients is used in this identity.  In particular $\beta<a_{\rm e}$.

On the closed form domain, square completion reads
\[
 q_a[x+y]=\ip{S_ax}{x}
       +\norm{T_a^{1/2}(y+T_a^{-1}Bx)}^2.
\]
The map $(x,y)\mapsto(x,y+T_a^{-1}Bx)$ is boundedly invertible and
preserves that domain.  It transfers the negative index to $S_a$;
also $\ker(\mathsf W_3-aI)$ is isomorphic to $\ker S_a$.
The negative even trial direction therefore makes the first even
Schur eigenvalue strictly negative, while its second is strictly
positive.  There is no even eigenvalue exactly at $a_{\rm e}$, and
exactly one even eigenvalue lies below it, counted with multiplicity.
Every odd eigenvalue exceeds $a_{\rm o}$.  The established compact
resolvent makes these full spectral statements.  Finally
Proposition~\ref{prop:v116-window-positive} supplies $\mu_0>0$,
and min--max applied to $v$ supplies $\mu_0\le\beta$.
Since $\beta<a_{\rm o}<a_{\rm e}$, the unique even low eigenvalue is
the simple ground state of the whole operator.  This proves
\eqref{eq:v117-ground-order} and the stronger even-sector threshold.
\end{proof}

\begin{corollary}[An ordinary ground-overlap bound at one window]
\label{cor:v117-ground-overlap}
For the normalized rational candidate $u$ of
Proposition~\ref{prop:v112-finite-certificate} and the complete unit
ground vector $\xi_0$ at $\lambda=3$,
\[
 \sin\angle(u,\xi_0)<0.01931.
\]
The same strict upper bound holds for the normalized exact $N=64$
source projection after the already certified normalization error is
included.  This is an ordinary Hilbert-space bound; it asserts no
endpoint comparison with $\xi_0$.
\end{corollary}
\begin{proof}
The improved trial $v$ in the proof above is even.  In its spectral
expansion, every non-ground component has energy greater than
$a_{\rm e}=10^{-34}$, and the ground energy is positive.  Hence
\[
 \sin\angle(v,\xi_0)
 <\sqrt{3.644\cdot10^{-38}/10^{-34}}.
\]
Also $z\perp u$, so the norm of the difference of the two rank-one
orthogonal projections onto $u$ and $v$ equals
$q/\sqrt{1+q^2}<0.000216$.  The triangle inequality for these
projectors proves
\[
 \sin\angle(u,\xi_0)
 <0.000216+\sqrt{3.644\cdot10^{-4}}
 <0.019305264<0.01931.
\]
For the unnormalized rational candidate $c$,
Proposition~\ref{prop:v113-source-certificate} gives
$\|P_{64}p_3-c\|<\varepsilon:=2.60\cdot10^{-36}$, while the
exact rational norm calculation gives $\|c\|>0.65172555$.
Writing $\rho=\varepsilon/\|c\|$, elementary ball geometry yields
\[
 \left\|\frac{P_{64}p_3}{\|P_{64}p_3\|}-\frac c{\|c\|}\right\|^2
 \le\frac{2\rho^2}{1+\sqrt{1-\rho^2}}<(4\cdot10^{-36})^2.
\]
Indeed the real inner product of these two unit vectors is positive
and at least $\sqrt{1-\rho^2}$, by minimizing the squared distance
from $c$ along the positive ray through $P_{64}p_3$.
This normalization error fits within the final margin.  The ordinary
norm controls projector distance but not the physical endpoint
functional; the earlier endpoint obstruction is therefore unchanged.
\end{proof}

This proves parity, simplicity and quantitative ordinary overlap at one
fixed window. It does not identify the unprojected source with the
ground vector. There is still no uniform family of spectral thresholds
or source overlap estimates along unbounded windows. Endpoint stability,
the existing sampler graph defect and full-strip transfer requirements
remain separate; G2 and RH are not closed by this result.

